\section{Conclusions}
\label{sec:conclusions}

This paper develops a risk-aware planning framework that connects
mission-dependent graph search with fixed-wing trajectory generation.
Hard no-fly constraints and continuous soft exposure have distinct roles
in the model, while time, exposure and fuel are combined through shared
reference scales. Coarse and corridor-constrained fine searches retain
the same mission objective. Subsequent spline refinement and speed
assignment convert the selected path into a trajectory evaluated against
the prescribed flight envelope. Reintegrating the mission metrics after
these stages preserves the distinction between search estimates and the
cost of the resulting motion.

%% EMPIRICAL CONCLUSIONS: 在此插入第四章支持的一个结果段，建议3--4句：
%% (a) 指明主案例权重和时间基线，填写C的变化及T/F实际代价；
%% (b) 指明同次main的粗细路径变化，不把样条/速度的作用归给fine搜索；
%% (c) 给出最终曲率、爬升角、速度/加速度和硬区检查的实际结果；
%% (d) 从72组总结一条本场景内的权重响应规律，注明其为coarse+后处理协议。
%% 不把“需要证明的问题”写成“已经证明的结论”，也不要求三个代价全部下降。

The scope of the study is a prescribed static environment and a reduced
aircraft model. The exposure integral is a planning surrogate, and the
fuel estimate is evaluated under the stated aerodynamic and efficiency
assumptions. Numerical feasibility checks do not constitute a continuous
collision certificate or a closed-loop flight demonstration. Within this
scope, the framework provides an auditable basis for relating mission
preferences to the geometry, motion limits and evaluated cost of a planned
trajectory. Further work will examine the effect of environmental and
model uncertainty and connect the generated references with higher-fidelity
flight dynamics and trajectory tracking.
